\section{Conclusions}

The critical deficit in continuous air quality monitoring across rapidly urbanizing Indian cities necessitates accurate virtual sensing frameworks. However, standard methodologies persistently suffer from peak-smoothing, often attenuating the magnitude of severe winter smog events. This study proposed a Two-Stage Decoupled Framework designed to combine daily spatial background modeling from Sentinel-5P satellite retrievals with high-frequency diurnal dynamics predicted by a 1D-CNN and BiLSTM network.

The proposed hybrid model improved the combined held-out-city $R^2$ from 0.4890 for the daily-baseline model to 0.6119, and reduced RMSE by 3.41 $\mu g/m^3$. In a severe evaluated event, the model generated a 6.725$\times$ hourly ratio and closely tracked the observed peak. These results indicate that decoupling daily spatial scale from hourly meteorological modulation is a promising approach for $PM_{2.5}$ virtual sensing. Future work must prioritize evaluating the framework using operational weather forecasts instead of historical reanalysis, conducting multi-fold spatial cross-validation, and integrating independent sensor networks for hyper-local validation.
